\section{Limitations}\label{sec:limitations}

Motivo deliberately targets offline, rule-based expansion into Standard MIDI Files.
That focus enables determinism and a compact implementation, but it also bounds what the system claims to replace.

\topic{Expressive and structural scope.}
The intermediate representation stores pitch, start beat, duration, and a fixed default velocity.
Per-note dynamics, continuous controllers, pitch bend, and synthesis parameters are not part of the source language.
Instruments are a closed keyword set mapped to General MIDI; compositions must live in a single file with no
\texttt{import} mechanism.
The lowerer caps output at 20,000 note events.

\topic{Workflow fit.}
Piano-roll editing is fast for exploratory scratch work but brittle for repetition; Motivo inverts that trade-off.
It rewards composers who think in rules, loops, and variation, while hand-placed micro-editing, mixing, and real-time
manipulation remain downstream in a DAW.

\topic{Evaluation scope.}
The evaluation reports compilation time, source-to-event expansion, and automated test coverage---not whether composers
prefer Motivo over a DAW or learn the syntax quickly.
Those questions would require a separate user study; this thesis does not claim them.
